% !TEX TS-program = LuaLaTeX+se
% !TEX encoding = UTF-8

%notes: tests \gabcneumesnippet: a bare neume (no staff lines, no clef) set
%notes: inline inside running text, alone and with NABC voice 1 (above) and
%notes: voice 2 (below).

\documentclass[11pt]{article}
\usepackage{fontspec}
\usepackage{geometry}
\geometry{a4paper,margin=1.5cm}
\usepackage[allowdeprecated=false]{gregoriotex}
\setmainfont[
    Path = ../../../fonts/ ,
    Extension = .otf ,
    UprightFont = *-Regular ,
    UprightFeatures = { SmallCapsFont = *SC-Regular } ,
    BoldFont = *-Bold ,
    BoldFeatures = { SmallCapsFont = *SC-Bold } ,
    ItalicFont = *-Italic ,
    ItalicFeatures = { SmallCapsFont = *SC-Italic } ,
    BoldItalicFont = *-BoldItalic ,
    BoldItalicFeatures = { SmallCapsFont = *SC-BoldItalic } ,
    Ligatures = TeX
]{Alegreya}

\listfiles
\begin{document}

\section*{Bare neume, no voice}
\noindent a torculus \gabcneumesnippet{fgf} in the middle of a
sentence, followed by more text to check that the paragraph did not break
and that the line continues normally after the neume.

\bigskip
\section*{Bare neume with NABC voice 1 (above staff)}
\noindent with St.\ Gall notation above it:
\gabcneumesnippet[1]{fgf|to} and a punctum \gabcneumesnippet[1]{f|pu}.

\bigskip
\section*{Bare neume with NABC voices 1 and 2 (above and below staff)}
\noindent with notation above and below:
\gabcneumesnippet[2]{fgf|to|ta} and \gabcneumesnippet[2]{f|pu|vi}.

\bigskip
\section*{Several neumes forming a word}
\noindent
\gabcneumesnippet{f}\gabcneumesnippet{g}\gabcneumesnippet{h}\gabcneumesnippet{g}
in sequence, and the same with NABC:
\gabcneumesnippet[1]{f|vi}\gabcneumesnippet[1]{g|pu}\gabcneumesnippet[1]{h|to}

\end{document}
